% Advanced Practice Examination: Game Theory and Information Economics
% Solutions only. Questions are in Micro1Q.tex.
% Compile with:
%   latexmk -pdf -outdir=aux Micro1Q_solutions.tex

\documentclass[12pt,a4paper]{article}

% Packages
\usepackage[utf8]{inputenc}
\usepackage[margin=1in]{geometry}
\usepackage{amsmath}
\usepackage{amssymb}
\usepackage{amsfonts}
\usepackage{bm}
\usepackage{mathtools}
\usepackage{booktabs}
\usepackage{tabularx}
\usepackage{enumitem}
\usepackage{hyperref}

% Enumerate style
\setlist[enumerate,1]{label=\arabic*., leftmargin=*}

\title{\textbf{Advanced Practice Examination:}\\
       \Large Game Theory and Information Economics\\[0.5em]
       \normalsize Comprehensive Answer Key}
\author{ }
\date{ }

\begin{document}

\maketitle

\section{Simultaneous Games and Dominance Solvability}

\subsection*{Solution to Problem 1.1}

\paragraph{Iterated Dominance (IESDS/IEWDS).}
\begin{itemize}
  \item \textbf{IESDS Round 1.} Strategy \(a_3\) is strictly dominated by \(a_1\):
  \[
    4 > 3,\quad 2 > 1,\quad 8 > 6,\quad 2 > 1.
  \]
  Eliminate \(a_3\).
  \item \textbf{IESDS Round 2.} In the reduced \(3\times 4\) matrix, check Player~B's strategies. For the remaining rows \((a_1, a_2, a_4)\), \(b_3\) is dominated by \(b_2\) for most rows (e.g.\ for \(a_1\) and \(a_4\)), but not strictly for all. We therefore also inspect Player~A again: \(a_2\) is weakly dominated by \(a_1\),
  \[
    4 > 1,\quad 2 = 2,\quad 8 > 0,\quad 2 > 1.
  \]
  \item \textbf{IEWDS Round.} Eliminating \(a_2\) and \(b_4\) (weakly dominated by \(b_2\)) reduces the game to a \(2\times 2\) between \(a_1, a_4\) and \(b_1, b_2\). The Nash equilibrium of this reduced game is
  \[
    (a_1, b_2)
  \]
  with payoffs \((2,2)\).
\end{itemize}

\paragraph{Pareto efficient profiles.}
The Pareto efficient strategy profiles in the original \(4\times 4\) matrix are
\[
  (a_1, b_3),\quad (a_4, b_1),\quad (a_2, b_3).
\]
Note in particular that \((a_1, b_3)\) gives Player~A their maximum payoff of \(8\).

\paragraph{``So What?'' evaluation.}
The Nash equilibrium \((a_1, b_2)\) yields payoffs \((2,2)\), while \((a_4, b_1)\) yields \((5,3)\). Both players strictly prefer \((a_4, b_1)\), but if Player~B believes Player~A will play \(a_4\), then deviating from \(b_1\) to \(b_2\) is profitable. Hence \((a_4, b_1)\) is not a Nash equilibrium. This mirrors a ``Fishermen's Dilemma'': individual incentives to ``overfish'' (deviate) drive the outcome toward a stable but Pareto-inefficient equilibrium, even though cooperation would yield higher payoffs.

\section{Sequential Interaction, Backward Induction, and Bargaining}

\subsection*{Solution to Problem 2.2 (Finite Horizon Bargaining)}

We solve by backward induction for \(T = 3\).
\begin{itemize}
  \item \textbf{Period \(T = 3\).} Player~1 makes the final offer. Since any positive amount is strictly preferred by Player~2 to receiving \(0\), Player~1 can offer
  \[
    (1,0),
  \]
  which is accepted.
  \item \textbf{Period \(T = 2\).} The surplus at the start of period \(2\) is \(\delta\). Player~2 must give Player~1 at least their continuation value from period \(3\), discounted back one period. At period \(3\), Player~1 would obtain \(\delta\), so at period \(2\) Player~2 must offer
  \[
    (\delta,\, 1-\delta).
  \]
  \item \textbf{Period \(T = 1\).} Player~1 must give Player~2 at least their discounted continuation value from period \(2\), which is \(\delta(1-\delta)\). Thus Player~1 proposes
  \[
    \bigl(1 - \delta + \delta^2,\; \delta - \delta^2\bigr),
  \]
  which Player~2 accepts in equilibrium.
\end{itemize}

As \(T \to \infty\), the equilibrium split converges to
\[
  \left(\frac{1}{1+\delta},\, \frac{\delta}{1+\delta}\right).
\]
Higher \(\delta\) (greater patience) reduces the first-mover advantage because the responder's threat to wait for a better future offer becomes more credible.

\section{Mixed Strategies, Security Payoffs, and Zero-Sum Dynamics}

\subsection*{Solution to Problem 3.1 (Military Attack)}

\paragraph{(i) Non-existence of pure-strategy Nash equilibrium.}
Check best responses:
\begin{itemize}
  \item If Player~1 plays \(A\), Player~2's best response is \(A\) (yielding \(0\) instead of \(4\)).
  \item If Player~2 plays \(A\), Player~1's best response is \(B\) (payoff \(3\) instead of \(0\) or \(2\)).
\end{itemize}
Iterating this logic over all cells shows that for any candidate pure profile, at least one player has a profitable deviation. Hence no pure-strategy Nash equilibrium exists.

\paragraph{(ii) Maxmin (security) strategy for Player 1.}
Let Player~1's mixed strategy be \((\sigma_1,\sigma_2,\sigma_3)\) over \((A,B,C)\). The expected payoffs if Player~2 defends each target are
\begin{align*}
  \mathbb{E}[U_1 \mid A] &= 0\cdot \sigma_1 + 3\sigma_2 + 2(1-\sigma_1-\sigma_2)
    = 2 - 2\sigma_1 + \sigma_2,\\
  \mathbb{E}[U_1 \mid B] &= 4\sigma_1 + 0\cdot \sigma_2 + 2(1-\sigma_1-\sigma_2)
    = 2 + 2\sigma_1 - 2\sigma_2,\\
  \mathbb{E}[U_1 \mid C] &= 4\sigma_1 + 3\sigma_2 + 0\cdot (1-\sigma_1-\sigma_2)
    = 4\sigma_1 + 3\sigma_2.
\end{align*}
In a security (maxmin) strategy, Player~1 chooses \((\sigma_1,\sigma_2,\sigma_3)\) so that these three expressions are equal, making Player~2 indifferent across pure strategies. Equating:
\begin{align*}
  2 - 2\sigma_1 + \sigma_2 &= 2 + 2\sigma_1 - 2\sigma_2
  \quad\Rightarrow\quad 4\sigma_1 = 3\sigma_2 \Rightarrow \sigma_2 = \tfrac{4}{3}\sigma_1,\\
  2 + 2\sigma_1 - 2\sigma_2 &= 4\sigma_1 + 3\sigma_2
  \quad\Rightarrow\quad
  2 = 2\sigma_1 + 5\sigma_2.
\end{align*}
Substituting \(\sigma_2 = \tfrac{4}{3}\sigma_1\) into the second equation gives
\[
  2 = 2\sigma_1 + 5\cdot \tfrac{4}{3}\sigma_1 = \tfrac{26}{3}\sigma_1
  \quad\Rightarrow\quad
  \sigma_1 = \frac{3}{13},\quad
  \sigma_2 = \frac{4}{13},\quad
  \sigma_3 = 1 - \sigma_1 - \sigma_2 = \frac{6}{13}.
\]

\paragraph{(iii) Value of the game.}
Substituting into, say, \(\mathbb{E}[U_1 \mid C]\),
\[
  v
  = 4\cdot \frac{3}{13} + 3\cdot \frac{4}{13}
  = \frac{12 + 12}{13}
  = \frac{24}{13}
  \approx 1.85.
\]
This is both the security level and the Nash equilibrium payoff for Player~1.

\paragraph{(iv) Comparison with an Inspection Game.}
In strictly zero-sum games, the Minmax Theorem implies the equality of the maxmin and minmax values and yields a saddle point: \(v_1 = -v_2\). In a non-zero-sum ``Inspection Game'', the employer's secure monitoring strategy and the worker's secure effort strategy do not generate such a saddle point because gains and losses are not perfectly opposed. As a result, the equality \(v_1 = -v_2\) fails, and equilibrium payoffs depend on the specific strategic interaction structure rather than a single game value.

\section{Adverse Selection and Market Signaling}

\subsection*{Solution to Problem 4.1}

\paragraph{(i) Optimal TIOLI price and trade probability.}
If the buyer offers a price \(p\), only sellers with \(\theta \leq p\) accept. Since \(\theta \sim \mathcal{U}[0,2]\),
\[
  \mathbb{E}[\theta \mid \theta \le p] = \frac{p}{2}.
\]
The buyer's expected utility is
\[
  \mathbb{E}[U_B]
  = \bigl(\mathbb{E}[\theta \mid \theta \le p] + a - p\bigr)\cdot \Pr(\theta \le p)
  = \left(\frac{p}{2} + a - p\right)\cdot \frac{p}{2}
  = \left(a - \frac{p}{2}\right)\frac{p}{2}.
\]
Maximising with respect to \(p\) yields the first-order condition
\[
  \frac{\partial}{\partial p}\left[\left(a - \frac{p}{2}\right)\frac{p}{2}\right]
  = \frac{a}{2} - \frac{p}{2} = 0
  \quad\Rightarrow\quad p^* = a.
\]
For all \(\theta \in [0,2]\) to trade, we need \(p^* \geq \max_{\theta}\theta = 2\), so
\[
  a \geq 2
\]
is the threshold ensuring a trade probability of \(1\).

\paragraph{(ii) Signaling via disclosure.}
If the seller reveals type \(\theta\) at cost \(c\), the buyer offers \(V_B(\theta) = \theta + a\), so the seller's payoff is
\[
  \theta + a - c.
\]
If the seller does not reveal, they receive \(p^* = a\). The seller reveals if
\[
  \theta + a - c \;\geq\; a
  \quad\Longleftrightarrow\quad
  \theta \geq c.
\]
Thus all types with \(\theta \geq c\) optimally disclose. As \(a\) increases, the buyer's valuation premium over the seller's own valuation grows, increasing the information rent that high-\(\theta\) sellers can capture (subject to the fixed disclosure cost \(c\) not being too large).

\section{Moral Hazard and Monopolistic Screening}

\subsection*{Solution to Problem 5.1}

\paragraph{(i) IC and PC constraints.}
If the worker exerts effort (Work), their expected utility is
\[
  U^{W} = w - 50 = 100 - 50 = 50.
\]
If the worker shirks (Not Work), then with probability \(1-m\) they receive the wage \(w\), and with probability \(m\) they are caught and receive \(0\). Since shirking has no effort cost, expected utility is
\[
  U^{N} = (1-m)w = (1-m)\cdot 100.
\]
The incentive compatibility (IC) constraint \(U^{W} \geq U^{N}\) becomes
\[
  100 - 50 \geq 100(1-m)
  \quad\Rightarrow\quad
  50 \geq 100 - 100m
  \quad\Rightarrow\quad
  m \geq 0.5.
\]
The participation constraint (PC) requires that the worker's utility from accepting the contract is at least their outside option \(U_0 = 0\):
\[
  100 - 50 = 50 \geq 0,
\]
which is satisfied.

\paragraph{(ii) Principal's optimal monitoring probability.}
Output is \(Y = 200\) if the worker exerts effort (which the IC constraint guarantees when \(m \ge 0.5\)), and \(0\) otherwise. Expected profit of the Principal when effort is induced is
\[
  \Pi = 200 - 100 - 10m = 100 - 10m.
\]
Subject to \(m \ge 0.5\), the Principal minimises the monitoring cost by choosing
\[
  m^* = 0.5,
\]
yielding
\[
  \Pi^{\text{SB}} = 200 - 100 - 10(0.5) = 95.
\]

\paragraph{(iii) Agency cost and comparison to first-best.}
In the first-best (effort observable, \(m = 0\)), the Principal pays \(w = 100\) and obtains output \(Y = 200\), so
\[
  \Pi^{\text{FB}} = 200 - 100 = 100.
\]
The agency cost is the loss in profit relative to the first-best:
\[
  \text{Agency Cost} = \Pi^{\text{FB}} - \Pi^{\text{SB}} = 100 - 95 = 5.
\]
This efficiency loss represents resources wasted on monitoring that would be unnecessary under symmetric information about effort.

\end{document}

